% Satzspiegel, Schrift und Pakete des gedruckten Handbuchs.
%
% Diese Datei ist handgepflegt. Der Buchkoerper unter kapitel/ stammt aus dem
% einmaligen Lauf von tools/convert-book.js, die Gestaltung hier nicht.

\usepackage{fontspec}
\usepackage[ngerman]{babel}
\usepackage{csquotes}
\usepackage{microtype}

% Libertinus statt Inter. Der Plan sah Inter vor, damit Buch und Website
% dasselbe Schriftbild bekommen - im Repo liegt Inter aber nur als woff2
% (common/public/css/inter/), und das kann fontspec nicht einbinden. Eine
% Serifenschrift traegt einen 100-Seiten-Fliesstext im Druck ohnehin besser als
% eine Oberflaechen-Grotesk. Falls Inter doch gewuenscht ist: woff2 mit fonttools
% nach otf wandeln, dann hier eintragen.
\setmainfont{Libertinus Serif}
\setsansfont{Libertinus Sans}
\setmonofont{Libertinus Mono}

\usepackage{amsmath}
\usepackage{graphicx}
\graphicspath{{media/}}

% Logiktafeln. booktabs setzt Linien mit Abstand statt Gitter - bei
% Wahrheitstabellen aus Nullen und Einsen ist das der Unterschied zwischen
% lesbar und Tapete.
\usepackage{booktabs}
\renewcommand{\arraystretch}{1.15}

\usepackage[
  a5paper,
  inner=20mm, outer=15mm, top=18mm, bottom=20mm,
]{geometry}

% 56 Stichwoerter, aus den XE-Feldern des Manuskripts uebernommen.
\usepackage{imakeidx}
\makeindex[title=Stichwortverzeichnis, intoc]

\usepackage[hidelinks]{hyperref}

% Abbildungen sitzen mitten im Text und sollen dort bleiben.
\setkomafont{caption}{\small}
\setkomafont{captionlabel}{\small\bfseries}
\setcapindent{0pt}
